\begin{trapbox}

	\begin{description}[style=nextline, leftmargin=2.8em, labelsep=0.5em, itemsep=0.35em]
		\item[\textbf{\textcolor{CTrap}{易错点一（量词未换）}}]
			学生对“存在”命题进行否定时，只将谓词改为否定形式，量词仍保留“存在”，例如将“存在 $x$ 使得 $p(x)$”错误地否定为“存在 $x$ 使得 $\neg p(x)$”。
		\item[\textbf{\textcolor{CTrap}{正确理解（量词必须换）：}}]
			否定存在量词命题必须同时将量词改为全称，即 $\neg(\exists x\in M,p(x))\iff \forall x\in M,\neg p(x)$，量词与谓词均需变化。
		\item[\textbf{\textcolor{CTrap}{错因分析（否定含义误）：}}]
			误以为否定仅在命题上加“不”，未理解存在量词的逻辑含义——否定“存在”等价于“所有都不”。
	\end{description}

	\medskip
	\begin{description}[style=nextline, leftmargin=2.8em, labelsep=0.5em, itemsep=0.35em]
		\item[\textbf{\textcolor{CTrap}{易错点二（谓词否定不全）}}]
			否定含有不等式的谓词时，只改变不等号方向但遗漏等号，例如将 $x+\frac{1}{x}>2$ 的否定错误写为 $x+\frac{1}{x}<2$。
		\item[\textbf{\textcolor{CTrap}{正确理解（大小等号要准确）：}}]
			“大于”的否定是“小于等于”，“小于”的否定是“大于等于”，“等于”的否定是“不等于”，必须包括所有相反情形。
		\item[\textbf{\textcolor{CTrap}{错因分析（思维定式忽略）：}}]
			对不等关系否定时只想到反向不等号，忽略了等号也成为可能性之一。
	\end{description}

	\medskip
	\begin{description}[style=nextline, leftmargin=2.8em, labelsep=0.5em, itemsep=0.35em]
		\item[\textbf{\textcolor{CTrap}{易错点三（误解集合变盘）}}]
			否定时认为集合 $M$ 也要变为补集，例如将“存在 $x\in M$”的否定写成“对任意 $x\notin M$”。
		\item[\textbf{\textcolor{CTrap}{正确理解（集合保持不变）：}}]
			否定只改变量词和谓词，不改变 $x$ 的取值范围 $M$，即 $\neg(\exists x\in M,p(x))\iff \forall x\in M,\neg p(x)$，$M$ 不变。
		\item[\textbf{\textcolor{CTrap}{错因分析（概念混淆误）：}}]
			将命题的否定与命题的逆否命题或集合的补集运算混淆。
	\end{description}

\end{trapbox}
